\section{Introduction}
\label{sec:introduction}

Long-context sequence modeling is constrained by three resources that cannot
be optimized independently: representational capacity, memory traffic, and
arithmetic work. Full causal self-attention preserves direct token-to-token
access, but its key-value state grows with the processed sequence and its
prefill attention matrix grows quadratically without additional structure.
Recurrent and linear-attention models replace that growing state with a
fixed-size summary, but any finite summary must eventually merge, overwrite,
or discard independent information. Sparse caches preserve selected records
exactly, yet a selection rule can retain the wrong records. Compressed
episodic representations extend the age horizon, but compression creates a
measurable approximation error.

\ALUCLU starts from the premise that these are not four implementations of one
memory abstraction. They are four different physical memory regimes:

\begin{table}[h]
\centering
\caption{The four memory regimes used by \ALUCLU. ``Exact'' is local to the
stored representation and never means unbounded exact recall.}
\label{tab:memory-regimes}
\begin{tabularx}{\textwidth}{@{}p{0.18\textwidth}p{0.21\textwidth}
  p{0.23\textwidth}X@{}}
\toprule
Regime & Primary role & Physical budget & Principal failure mode \\
\midrule
Dense recurrence & Continuous global statistics & Fixed fast-weight matrix
per layer & Interference and decay \\
Local softmax & Precise recent interactions & Last \(W\) projected keys and
values & Hard recency horizon \\
Surprise cache & Sparse exceptional records & Highest-priority \(C_x\)
projected KV records & Selection mismatch and cache freezing \\
Episodic capsules & Longer temporal organization & Current segment, \(K\)
live segments, \(C\) rank-\(r\) archive capsules & Compression and overwrite \\
\bottomrule
\end{tabularx}
\end{table}

The architecture places these regimes in parallel inside selected transformer-
like residual blocks. Every layer has a Gated Delta Rule-2 recurrent lane.
Local, exact-cache, and episodic lanes are scheduled at configurable layer
intervals. They all consume the same normalized token representation rather
than forming a serial chain. Their outputs are independently normalized and
combined with learned token-wise coefficients. This arrangement gives the
model an explicit choice between a dense recurrent estimate, exact recent
evidence, a sparse set of high-write-residual events, and temporally segmented
compressed evidence.

\subsection{Design question}

The project asks a deliberately narrower and more testable question than
``What is the universally best language-model architecture?'':

\begin{quote}
Can one build an executable causal reference architecture whose persistent
decode state is bounded by configuration, whose heterogeneous memory losses
are visible in state and metrics, and whose interfaces permit exact
token-by-token auditing before specialized kernels or large-scale training?
\end{quote}

This question matters because algorithmic asymptotics alone do not determine
real performance. A theoretically constant-state recurrence can be slower than
optimized attention if implemented as a Python token loop. A low-rank archive
can have a long nominal horizon but poor task recall. A small cache can be
``exact'' for accepted keys while missing the one key a future query needs.
Consequently, \ALUCLU treats equations, code semantics, state bytes, benchmark
protocols, and claim boundaries as parts of one engineering object.

\subsection{Contributions}

The present release makes the following concrete contributions.

\begin{enumerate}
  \item It specifies and implements a four-lane decoder block with explicit
  scan, chunk, and recurrent-step interfaces and immutable external state
  containers.

  \item It couples a Gated Delta Rule-2 write residual to a fixed-capacity
  exact KV cache. The cache retains the highest observed scalar residual
  priorities under a deterministic strict-replacement policy.

  \item It implements a bounded episodic hierarchy with exact current/live
  factor sums, low-rank archive capsules, exact positive denominator
  statistics, a TensorSketch-assisted router, and a Frobenius error ledger.

  \item It derives the persistent payload, retained temporal budget, and
  token-cost terms from the actual state fields instead of using an informal
  ``constant memory'' label.

  \item It provides a mass-normalized routing construction that preserves the
  unweighted denominator mass in exact arithmetic when numerical safety
  clamps are inactive, while making the clamp qualification explicit.

  \item It supplies an MIT-licensed Python package, command-line tools,
  protocol manifests, multilingual documentation, reproducible JSON evidence,
  and an automated release gate with 50 tests.

  \item It defines a staged evaluation program that separates small
  correctness and learnability diagnostics from model-quality, accelerator,
  energy, and long-context claims.
\end{enumerate}

\subsection{Evidence contract}
\label{sec:evidence-contract-intro}

The strongest conclusions supported by the current artifact are mechanical:
the code executes, its configured state is bounded, tested causal paths agree
between full scans and recurrent stepping in evaluation mode, and the recorded
small CPU experiments are reproducible under their declared protocols. The
artifact does \emph{not} establish frontier language-model quality, universal
speed superiority, an optimal memory allocation, or a production GPU kernel.

We use three labels throughout the paper:

\begin{description}
  \item[Implemented] The behavior exists in the released source code and is
  connected to the public model path.
  \item[Measured] A result file records the protocol, configuration,
  environment, and numerical outcome.
  \item[Proposed] The experiment or optimization is fully specified but has
  not been executed as evidence for this release.
\end{description}

This vocabulary is intentionally conservative. It prevents a common category
error in architecture research: presenting a promising mechanism, a Python
prototype, or an asymptotic argument as though it were already an end-to-end
systems result.

\subsection{Paper organization}

\Cref{sec:related-work} positions the design among recurrent, attention,
hybrid, and long-context evaluation work. \Cref{sec:problem-contract} defines
the causal interface and measurable Pareto contract. \Cref{sec:architecture}
describes the block and execution order. \Cref{sec:memory-lanes} gives the
implemented equations, while \Cref{sec:theory} derives state, complexity, and
error properties. \Cref{sec:implementation} maps the specification to code.
\Cref{sec:evaluation,sec:results} separate the planned evaluation matrix from
the measurements already available. \Cref{sec:limitations} lists failure
modes and physical limits. The appendices provide configuration, state, test,
and reproduction details.
